% ══════════════════════════════════════════════════════════════════
\chapter{Introduction and Mathematical Framework}
\label{ch:framework}
% ══════════════════════════════════════════════════════════════════

\section{Problem Statement}

Consider a deep feed-forward ReLU network with $L$ hidden layers.
Layer~$l$ maps $\R^{d_{l-1}} \to \R^{d_l}$ via
\begin{equation}\label{eq:forward}
  \bm{h}^{(0)} = \bm{x}, \qquad
  \bm{h}^{(l)} = \relu\!\bigl(W^{(l)}\,\bm{h}^{(l-1)}\bigr),
  \quad l = 1,\dots,L,
\end{equation}
where $W^{(l)} \in \R^{d_l \times d_{l-1}}$ and
$\relu(z) = \max(0,z)$ is applied element-wise.
In general, the weight matrices need not be square.
For simplicity, in this work we study constant-width networks
($d_l = d$ for all~$l$), but the analysis generalizes to
non-square layers.

We study two desiderata for the weight matrices $\{W^{(l)}\}$:
\begin{enumerate}[label=(\roman*)]
  \item \textbf{Geometry preservation.}  Class separability of the input
        data should survive the composition of $L$ layers.  Formally,
        for data pairs $(\bm{x}_i, \bm{x}_j)$, we would like
        \[
          \E\!\Bigl[
            \bigl\langle \relu(W\bm{x}_i),\; \relu(W\bm{x}_j) \bigr\rangle
          \Bigr]
          \;\approx\;
          \bigl\langle \bm{x}_i,\, \bm{x}_j \bigr\rangle.
        \]
  \item \textbf{Gradient stability.}  The gradient magnitude
        $\|\partial \mathcal{C}/\partial W^{(l)}\|$ should be approximately
        the same for every layer~$l$, so that all layers receive a useful
        learning signal.
\end{enumerate}

\noindent
The central finding of this work is that these two objectives are in
\emph{structural conflict} when the nonlinearity is ReLU.

% ──────────────────────────────────────────────────────────────────
\section{The Arc-Cosine Kernel (Summary)}
\label{sec:kernel}

The arc-cosine kernel framework~\cite{cho2009kernel} characterizes
the expected inner product of post-ReLU outputs.
These results were discussed previously with the advisor; we state
them here without proof for reference.

Let $W \in \R^{d_{\mathrm{out}} \times d_{\mathrm{in}}}$ have i.i.d.\
rows $\bm{w}_k \sim \mathcal{N}(\bm{0},\, \sigma^2 I)$.
For two unit-norm inputs with angle
$\alpha = \arccos\!\bigl(\langle \bm{x}_i, \bm{x}_j \rangle\bigr)
  \in [0,\pi]$:

\begin{theorem}[Arc-cosine kernel, first order]
\label{thm:kernel}
\begin{equation}\label{eq:kernel}
  K(\alpha)
  \;=\;
  \frac{1}{d_{\mathrm{out}}}
  \,\E\!\Bigl[
    \bigl\langle \relu(W\bm{x}_i),\; \relu(W\bm{x}_j) \bigr\rangle
  \Bigr]
  \;=\;
  \frac{\sigma^2}{2\pi}\,
  \bigl(\sin\alpha + (\pi - \alpha)\cos\alpha\bigr).
\end{equation}
\end{theorem}

The normalized kernel $\hat{K}(\alpha) = K(\alpha)/K(0)
= \frac{1}{\pi}(\sin\alpha + (\pi-\alpha)\cos\alpha)$
gives the expected cosine similarity of the outputs.
The output angle after one layer is:
\begin{equation}\label{eq:angle_map}
  \alpha' = \arccos\!\bigl(\hat{K}(\alpha)\bigr)
  = \arccos\!\left(
    \frac{\sin\alpha + (\pi-\alpha)\cos\alpha}{\pi}
  \right).
\end{equation}
Since $\hat{K}(\alpha) > \cos\alpha$ for all $\alpha \in (0,\pi)$,
the output angle is always smaller than the input angle:

\begin{proposition}[Systematic angle contraction]
\label{prop:distortion}
\begin{equation}\label{eq:distortion}
  D(\alpha)
  = \hat{K}(\alpha) - \cos\alpha
  = \frac{1}{\pi}\bigl(\sin\alpha - \alpha\cos\alpha\bigr)
  > 0
  \quad\text{for all } \alpha \in (0, \pi).
\end{equation}
\end{proposition}

\noindent
This means \textbf{every pair of distinct vectors becomes more similar
after one ReLU layer}.  Over $L$ layers this compounds: angles contract
toward~$0$, causing \emph{geometric collapse}.
This distortion is independent of $\sigma^2$ and is therefore
a property of the ReLU nonlinearity, not of the initialization.
The empirical consequences are explored in
Chapters~\ref{ch:geometry} and~\ref{ch:angle_map}.

% ──────────────────────────────────────────────────────────────────
\section{Measuring Geometry}
\label{sec:geometry_metrics}

Previous work in this project used a PCA-based metric:
$\Var(\text{PC}_1)/\bigl(\Var(\text{PC}_1) + \Var(\text{PC}_2)\bigr)$,
which measures the elongation of the 2D PCA projection.
As we will show in Chapter~\ref{ch:geometry}, this metric is
\textbf{misleading}: data can be spread across many dimensions
(appearing ``well-preserved'' in PCA scatter plots) while having
\emph{no class structure} whatsoever.

We adopt three metrics that capture complementary aspects of geometry:

\begin{definition}[$k$-NN accuracy]
\label{def:knn}
Given projected data $\{\bm{h}_i^{(L)}\}_{i=1}^{n}$ with labels
$\{y_i\}$, the $k$-NN accuracy is the \emph{leave-one-out}
classification accuracy using Euclidean distance with $k=5$ neighbors.

Concretely: for each data point~$i$, we find the $k=5$ nearest
points in the projected space (by Euclidean distance among all
\emph{other} points), and classify~$i$ by majority vote of their
labels.  The $k$-NN accuracy is the fraction of points correctly
classified this way.

Chance level is $1/C$ where $C$ is the number of classes
(0.10 for Fashion-MNIST with 10 classes).
\end{definition}

\begin{remark}[Sensitivity to the choice of $k$]
The choice $k=5$ is not critical.  At $k=1$, the metric is more
sensitive to noise (a single nearest neighbor could be an outlier).
At $k=100$, it over-smooths and measures only coarse cluster
separation.  For our 2000-sample, 10-class dataset, $k=5$
provides a stable, local measure of class structure.
Empirically, results are similar for $k \in \{3, 5, 10\}$.
\end{remark}

\begin{definition}[Pairwise distance correlation]
Sample $m$ random pairs.  For each pair $(i,j)$, compute the
Euclidean distance $d^{(0)}_{ij} = \|\bm{x}_i - \bm{x}_j\|$
in the input space and
$d^{(L)}_{ij} = \|\bm{h}_i^{(L)} - \bm{h}_j^{(L)}\|$
in the projected space.  The metric is the Spearman rank correlation
$\rho_s\!\bigl(\{d^{(0)}_{ij}\}, \{d^{(L)}_{ij}\}\bigr)$.
\end{definition}

\begin{definition}[Effective dimensionality]
Let $\lambda_1 \geq \lambda_2 \geq \cdots \geq \lambda_d \geq 0$ be
the eigenvalues of the covariance matrix of $\{\bm{h}_i^{(L)}\}$.
Define the normalized distribution
$p_k = \lambda_k / \sum_j \lambda_j$.  The effective dimensionality is
\begin{equation}\label{eq:effdim}
  d_{\mathrm{eff}}
  = \exp\!\Bigl(-\sum_{k} p_k \ln p_k\Bigr)
  = \exp\!\bigl(H(\{p_k\})\bigr),
\end{equation}
where $H$ is the Shannon entropy.  This ranges from~$1$ (all variance
in one direction) to~$d$ (uniform spread).
\end{definition}

\begin{remark}[Why $k$-NN is our primary metric]
\label{rem:knn_primary}
The $k$-NN accuracy directly answers the question:
\emph{can we tell which class a point belongs to based on its
neighbors in the projected space?}  This is the operational definition
of ``class structure survives the projection.''

Unlike distance correlation (which measures global distance ranking)
or effective dimensionality (which measures spread), $k$-NN measures
whether \emph{local neighborhoods} are class-consistent --- the
most relevant notion for classification.

In particular, high effective dimensionality combined with low
$k$-NN accuracy indicates \emph{uniformly dispersed noise}
(data spread across many dimensions with no class structure),
as we will see for row-centered initialization in
Chapter~\ref{ch:geometry}.
\end{remark}

% ──────────────────────────────────────────────────────────────────
\section{Measuring Gradient Health}
\label{sec:gradient_health}

\subsection{Activation and error signal matrices}

\begin{definition}[Activation matrix]
\label{def:activation_matrix}
Given a batch of $n$ input samples $\{\bm{x}_1, \dots, \bm{x}_n\}$,
the \emph{activation matrix} at layer~$l$ is
$A^{(l)} \in \R^{n \times d_l}$, where the $i$-th row is the
post-ReLU activation vector of the $i$-th sample:
\[
  (A^{(l)})_{ik} \;=\; a^{(l)}_{ik}
  \;=\; [\bm{h}^{(l)}_i]_k
  \;=\; \relu\!\bigl(z^{(l)}_{ik}\bigr),
\]
with pre-activation
$z^{(l)}_{ik} = \sum_{j=1}^{d_{l-1}} W^{(l)}_{kj}\, a^{(l-1)}_{ij}$.
In words: $a^{(l)}_{ik}$ is the value of the $k$-th neuron in
layer~$l$ when processing sample~$i$.
\end{definition}

\begin{definition}[Error signal matrix]
\label{def:error_signal}
For a loss function $\mathcal{C}$, the \emph{error signal}
(backpropagated error) at layer~$l$ for sample~$i$ is
\[
  \bm{\delta}^{(l)}_i
  \;=\; \frac{\partial \mathcal{C}}{\partial \bm{z}^{(l)}_i}
  \;\in\; \R^{d_l},
\]
where $\bm{z}^{(l)}_i$ is the pre-activation vector at layer~$l$
for sample~$i$.  This is computed recursively via the
\emph{backpropagation equation}:
\begin{equation}\label{eq:bp_recursion}
  \bm{\delta}^{(l)}_i
  = \bigl(W^{(l+1)}\bigr)^\top \bm{\delta}^{(l+1)}_i
    \;\odot\; \sigma'\!\bigl(\bm{z}^{(l)}_i\bigr),
\end{equation}
where $\sigma'(\bm{z}^{(l)}_i) = \mathbf{1}[\bm{z}^{(l)}_i > 0]$
is the ReLU derivative (a binary mask: $1$ where the pre-activation
was positive, $0$ where it was not).

The \emph{error signal matrix} $\Delta^{(l)} \in \R^{n \times d_l}$
collects these row-wise:
$(\Delta^{(l)})_{ik} = \delta^{(l)}_{ik}$.
\end{definition}

\subsection{Root-mean-square (RMS) signal levels}

Using the activation matrix from Definition~\ref{def:activation_matrix},
we define:
\begin{equation}\label{eq:rms}
  \rms\!\bigl(A^{(l)}\bigr)
  = \frac{\|A^{(l)}\|_F}{\sqrt{n \cdot d_l}}
  = \sqrt{\frac{1}{n \cdot d_l}
    \sum_{i=1}^{n}\sum_{k=1}^{d_l} (a^{(l)}_{ik})^2}
  = \sqrt{\E_{i,k}\!\bigl[(a^{(l)}_{ik})^2\bigr]}.
\end{equation}
This is the typical magnitude of a single activation entry, independent
of batch size and layer width.  The same definition applies to the
error signal matrix: $\rms(\Delta^{(l)})$.

\subsection{Per-layer gains}

\begin{definition}[Forward gain]
\begin{equation}\label{eq:fwd_gain}
  g_{\mathrm{fwd}}^{(l)}
  = \frac{\rms(A^{(l)})}{\rms(A^{(l-1)})}
  = \sqrt{\frac{\E[(a^{(l)})^2]}{\E[(a^{(l-1)})^2]}}.
\end{equation}
\end{definition}

\begin{definition}[Backward gain]
\label{def:bwd_gain}
\begin{equation}\label{eq:bwd_gain}
  g_{\mathrm{bwd}}^{(l)}
  = \frac{\rms(\Delta^{(l)})}{\rms(\Delta^{(l+1)})},
\end{equation}
using the error signal matrices from
Definition~\ref{def:error_signal}.
Note the direction: the backward gain measures how the error
signal grows (or shrinks) as it propagates \emph{backward}
from layer~$l+1$ to layer~$l$.
\end{definition}

\subsection{Connection to weight gradients}
\label{sec:connection_grad}

The fourth backpropagation equation gives the gradient of the loss
with respect to the weights at layer~$l$.  For a single sample~$i$:
\[
  \frac{\partial \mathcal{C}_i}{\partial W^{(l)}}
  = \bm{\delta}^{(l)}_i \,\bigl(\bm{a}^{(l-1)}_i\bigr)^\top
  \;\in\; \R^{d_l \times d_{l-1}}.
\]
Averaging over the batch of $n$ samples and writing in matrix form:
\begin{equation}\label{eq:bp4}
  \frac{\partial \mathcal{C}}{\partial W^{(l)}}
  = \frac{1}{n}\,
    \bigl(\Delta^{(l)}\bigr)^\top A^{(l-1)}
  \;\in\; \R^{d_l \times d_{l-1}}.
\end{equation}

We can now bound the gradient magnitude.  Taking the Frobenius norm:
\[
  \left\|\frac{\partial \mathcal{C}}{\partial W^{(l)}}\right\|_F
  = \frac{1}{n}\,
    \bigl\|(\Delta^{(l)})^\top A^{(l-1)}\bigr\|_F.
\]
By the submultiplicativity of the Frobenius norm
($\|AB\|_F \leq \|A\|_F\,\|B\|_F$):
\[
  \left\|\frac{\partial \mathcal{C}}{\partial W^{(l)}}\right\|_F
  \;\leq\;
  \frac{1}{n}\,
  \|\Delta^{(l)}\|_F \cdot \|A^{(l-1)}\|_F.
\]
Writing $\|\Delta^{(l)}\|_F = \rms(\Delta^{(l)}) \cdot \sqrt{n \cdot d_l}$
and $\|A^{(l-1)}\|_F = \rms(A^{(l-1)}) \cdot \sqrt{n \cdot d_{l-1}}$:
\begin{equation}\label{eq:grad_magnitude}
  \left\|\frac{\partial \mathcal{C}}{\partial W^{(l)}}\right\|_F
  \;\leq\;
  \sqrt{d_l \cdot d_{l-1}}\;\cdot\;
  \rms\!\bigl(\Delta^{(l)}\bigr) \;\cdot\;
  \rms\!\bigl(A^{(l-1)}\bigr).
\end{equation}
When the rows of $\Delta^{(l)}$ and $A^{(l-1)}$ are approximately
uncorrelated (which holds at initialization), this bound is
approximately tight.  Since $\sqrt{d_l \cdot d_{l-1}}$ is the same
for every layer (in a constant-width network), we conclude:
\begin{equation}\label{eq:grad_proportional}
  \left\|\frac{\partial \mathcal{C}}{\partial W^{(l)}}\right\|_F
  \;\propto\;
  \rms\!\bigl(A^{(l-1)}\bigr) \;\cdot\;
  \rms\!\bigl(\Delta^{(l)}\bigr).
\end{equation}
The gradient magnitude at layer~$l$ is proportional to the product
of the \textbf{forward signal} (activation RMS at layer~$l-1$) and
the \textbf{backward signal} (error signal RMS at layer~$l$).
If either varies across layers, the gradients become non-uniform.

\subsection{Why gain $\approx 1$ is critical}

If the per-layer gain is a constant $g$ across all $L$ layers, the
cumulative effect is $g^L$:
\begin{center}
\begin{tabular}{ccc}
\toprule
Per-layer gain $g$ & $g^{50}$ & Effect \\
\midrule
1.01 & 1.64 & Mild growth \\
1.10 & $1.17 \times 10^2$ & Explosion \\
\textbf{0.826} & $\bm{7.1 \times 10^{-5}}$ & \textbf{Vanishing} \\
0.99 & 0.61 & Mild decay \\
\bottomrule
\end{tabular}
\end{center}
The value $g = 0.826 = \sqrt{(\pi-1)/\pi}$ will appear throughout this
work as the structural forward gain of row-centered initialization
(Chapter~\ref{ch:gradient_trap}).
